% ============================================================================
% COMPREHENSIVE LATEX TEMPLATE COLLECTION
% Best Practices for Tables, Boxes, Figures, Workflows, and More
% ============================================================================

\documentclass[a4paper,12pt]{article}

% ============================================================================
% ESSENTIAL PACKAGES
% ============================================================================
\usepackage[utf8]{inputenc}
\usepackage[T1]{fontenc}
\usepackage{geometry}
\usepackage{graphicx}
\usepackage{xcolor}
\usepackage{colortbl}
\usepackage{booktabs}
\usepackage{array}
\usepackage{tabularx}
\usepackage{longtable}
\usepackage{tabulary}
\usepackage{multirow}
\usepackage{multicol}
\usepackage{siunitx}
\usepackage{tikz}
\usetikzlibrary{
    shapes,
    arrows,
    positioning,
    calc,
    trees,
    graphs,
    graphdrawing,
    shapes.geometric,
    shapes.arrows,
    shapes.multipart,
    chains,
    fit,
    backgrounds,
    shadows,
    decorations.pathreplacing,
    decorations.markings,
    patterns,
    automata,
    matrix,
    intersections
}
\usegdlibrary{trees, layered, force}
\usepackage{tcolorbox}
\tcbuselibrary{
    skins,
    breakable,
    theorems,
    listings,
    xparse,
    hooks,
    raster
}
\usepackage{listings}
\usepackage{minted}
\usepackage{float}
\usepackage{caption}
\usepackage{subcaption}
\usepackage{wrapfig}
\usepackage{framed}
\usepackage{mdframed}
\usepackage{fancybox}
\usepackage{tikz-cd}
\usepackage{smartdiagram}
\usepackage{chemfig}
\usepackage{pgfgantt}
\usepackage{pgffor}
\usepackage{amsmath}
\usepackage{amssymb}
\usepackage{amsfonts}
\usepackage{mathtools}
\usepackage{nccmath}
\usepackage{hyperref}
\usepackage{cleveref}
\usepackage{enumitem}
\usepackage{parskip}
\usepackage{titlesec}
\usepackage{fancyhdr}
\usepackage{lastpage}
\usepackage{draftwatermark}
\usepackage{background}
\usepackage{eso-pic}
\usepackage{lipsum} % For dummy text

% ============================================================================
% PAGE SETUP
% ============================================================================
\geometry{
    top=2.5cm,
    bottom=2.5cm,
    left=2.5cm,
    right=2.5cm,
    headheight=15pt
}

% ============================================================================
% COLOR DEFINITIONS
% ============================================================================
\definecolor{primary}{RGB}{0, 51, 102}
\definecolor{secondary}{RGB}{0, 102, 153}
\definecolor{accent}{RGB}{255, 102, 0}
\definecolor{lightgray}{RGB}{240, 240, 240}
\definecolor{darkgray}{RGB}{80, 80, 80}
\definecolor{success}{RGB}{0, 153, 0}
\definecolor{warning}{RGB}{255, 153, 0}
\definecolor{danger}{RGB}{204, 0, 0}
\definecolor{info}{RGB}{0, 153, 204}

% ============================================================================
% CODE LISTINGS SETUP
% ============================================================================
\definecolor{codegreen}{rgb}{0,0.6,0}
\definecolor{codegray}{rgb}{0.5,0.5,0.5}
\definecolor{codepurple}{rgb}{0.58,0,0.82}
\definecolor{backcolour}{rgb}{0.95,0.95,0.92}

\lstdefinestyle{mystyle}{
    backgroundcolor=\color{backcolour},
    commentstyle=\color{codegreen},
    keywordstyle=\color{magenta},
    numberstyle=\tiny\color{codegray},
    stringstyle=\color{codepurple},
    basicstyle=\ttfamily\footnotesize,
    breakatwhitespace=false,
    breaklines=true,
    captionpos=b,
    keepspaces=true,
    numbers=left,
    numbersep=5pt,
    showspaces=false,
    showstringspaces=false,
    showtabs=false,
    tabsize=2,
    frame=single,
    framerule=0.5pt,
    rulecolor=\color{codegray}
}

\lstset{style=mystyle}

% ============================================================================
% TCOLORBOX DEFINITIONS
% ============================================================================
% Basic colored box
\newtcolorbox{basicbox}[1][]{
    colback=blue!5!white,
    colframe=blue!75!black,
    fonttitle=\bfseries,
    #1
}

% Theorem-style box
\newtcolorbox{theobox}[2][]{
    colback=green!5!white,
    colframe=green!75!black,
    fonttitle=\bfseries,
    title=#2,
    theorem style=standard,
    #1
}

% Example box with shadow
\newtcolorbox{examplebox}[2][]{
    enhanced,
    colback=yellow!10!white,
    colframe=yellow!75!black,
    fonttitle=\bfseries,
    title=#2,
    drop shadow,
    #1
}

% Code listing box
\newtcblisting{codebox}[2][]{
    enhanced,
    colback=white,
    colframe=gray!75!black,
    fonttitle=\bfseries,
    title=#2,
    listing only,
    listing options={style=mystyle},
    left=6pt,
    right=6pt,
    top=6pt,
    bottom=6pt,
    boxrule=0.5pt,
    #1
}

% Breakable box for long content
\newtcolorbox{breakablebox}[2][]{
    enhanced jigsaw,
    breakable,
    colback=white,
    colframe=primary,
    fonttitle=\bfseries,
    title=#2,
    sharp corners,
    rounded corners=southeast,
    arc is angular,
    arc=3mm,
    underlay first={
        \path[tcolorframe] ([xshift=2mm]frame.north west) -- ([xshift=2mm]frame.north east);
    },
    underlay middle={
        \path[tcolorframe] ([xshift=2mm]frame.north west) -- ([xshift=2mm]frame.north east);
        \path[tcolorframe] ([xshift=2mm]frame.south west) -- ([xshift=2mm]frame.south east);
    },
    underlay last={
        \path[tcolorframe] ([xshift=2mm]frame.south west) -- ([xshift=2mm]frame.south east);
    },
    #1
}

% Info/Warning/Danger boxes
\newtcolorbox{infobox}[2][]{
    colback=info!10!white,
    colframe=info,
    fonttitle=\bfseries,
    title=#2,
    attach boxed title to top left={yshift=-2mm, xshift=2mm},
    boxed title style={
        colback=info,
        colframe=info,
        arc=2mm
    },
    #1
}

\newtcolorbox{warningbox}[2][]{
    colback=warning!10!white,
    colframe=warning,
    fonttitle=\bfseries,
    title=#2,
    attach boxed title to top left={yshift=-2mm, xshift=2mm},
    boxed title style={
        colback=warning,
        colframe=warning,
        arc=2mm
    },
    #1
}

\newtcolorbox{dangerbox}[2][]{
    colback=danger!10!white,
    colframe=danger,
    fonttitle=\bfseries,
    title=#2,
    attach boxed title to top left={yshift=-2mm, xshift=2mm},
    boxed title style={
        colback=danger,
        colframe=danger,
        arc=2mm
    },
    #1
}

% Side-by-side boxes (raster)
\newtcbinputlisting{\syntaxbox}[3][]{
    listing only,
    listing options={#2},
    title=#3,
    colback=white,
    colframe=darkgray,
    fonttitle=\bfseries,
    #1
}

% ============================================================================
% TABLE TEMPLATES
% ============================================================================
% Professional table with booktabs
\newcolumntype{L}[1]{>{\raggedright\let\newline\\\arraybackslash\hspace{0pt}}m{#1}}
\newcolumntype{C}[1]{>{\centering\let\newline\\\arraybackslash\hspace{0pt}}m{#1}}
\newcolumntype{R}[1]{>{\raggedleft\let\newline\\\arraybackslash\hspace{0pt}}m{#1}}

% Colored row command
\newcommand{\hrowcolor}[1]{\rowcolor{#1}}

% Alternating row colors
\usepackage{etoolbox}
\newcounter{tablerow}
\newcommand{\alternaterow}{%
    \stepcounter{tablerow}%
    \ifodd\value{tablerow}%
        \rowcolor{lightgray}%
    \fi%
}

% ============================================================================
% FIGURE AND IMAGE SETUP
% ============================================================================
\captionsetup{
    font=small,
    labelfont=bf,
    labelsep=colon,
    justification=justified,
    singlelinecheck=false
}

% ============================================================================
% TIKZ STYLES FOR DIAGRAMS
% ============================================================================
\tikzset{
    % Basic node styles
    block/.style={
        rectangle,
        draw,
        fill=blue!10,
        text width=6em,
        text centered,
        rounded corners,
        minimum height=3em,
        drop shadow
    },
    line/.style={
        draw,
        -latex',
        thick
    },
    cloud/.style={
        draw,
        ellipse,
        fill=red!10,
        minimum height=2em,
        drop shadow
    },
    decision/.style={
        diamond,
        draw,
        fill=green!10,
        text width=4.5em,
        text centered,
        inner sep=0pt,
        drop shadow
    },
    startstop/.style={
        rectangle,
        rounded corners,
        minimum width=3cm,
        minimum height=1cm,
        text centered,
        draw=black,
        fill=red!10,
        drop shadow
    },
    process/.style={
        rectangle,
        minimum width=3cm,
        minimum height=1cm,
        text centered,
        draw=black,
        fill=blue!10,
        drop shadow
    },
    io/.style={
        trapezium,
        trapezium left angle=70,
        trapezium right angle=110,
        minimum width=3cm,
        minimum height=1cm,
        text centered,
        draw=black,
        fill=green!10,
        drop shadow
    },
    arrow/.style={
        thick,
        ->,
        >=stealth
    },
    dashedarrow/.style={
        thick,
        ->,
        >=stealth,
        dashed
    },
    bidirectional/.style={
        thick,
        <->,
        >=stealth
    }
}

% ============================================================================
% DOCUMENT METADATA
% ============================================================================
\title{\textbf{Comprehensive LaTeX Template Collection}\\
\large Best Practices for Tables, Boxes, Figures, Workflows \& Schemas}
\author{LaTeX Template Repository}
\date{\today}

% ============================================================================
% BEGIN DOCUMENT
% ============================================================================
\begin{document}

\maketitle

\begin{abstract}
This document provides a comprehensive collection of LaTeX templates and best 
practices for creating professional documents. It includes examples of tables, 
colored boxes (tcolorbox), figures, workflows, schemas, code listings, and more.
All templates are ready to use and can be customized for your specific needs.
\end{abstract}

\tableofcontents
\newpage

% ============================================================================
% SECTION 1: TABLES
% ============================================================================
\section{Professional Table Templates}

\subsection{Basic Professional Table with Booktabs}
\begin{table}[H]
\centering
\caption{Simple Professional Table}
\label{tab:simple}
\begin{tabular}{lcc}
\toprule
\textbf{Column 1} & \textbf{Column 2} & \textbf{Column 3} \\
\midrule
Row 1, Col 1 & Row 1, Col 2 & Row 1, Col 3 \\
Row 2, Col 1 & Row 2, Col 2 & Row 2, Col 3 \\
Row 3, Col 1 & Row 3, Col 2 & Row 3, Col 3 \\
\bottomrule
\end{tabular}
\end{table}

\subsection{Table with Colored Rows}
\begin{table}[H]
\centering
\caption{Table with Alternating Row Colors}
\label{tab:colored}
\begin{tabular}{lcc}
\toprule
\rowcolor{primary!20}
\textbf{Header 1} & \textbf{Header 2} & \textbf{Header 3} \\
\midrule
\rowcolor{lightgray}
Data 1 & Data 2 & Data 3 \\
Data 4 & Data 5 & Data 6 \\
\rowcolor{lightgray}
Data 7 & Data 8 & Data 9 \\
Data 10 & Data 11 & Data 12 \\
\bottomrule
\end{tabular}
\end{table}

\subsection{Table with Multi-row and Multi-column}
\begin{table}[H]
\centering
\caption{Complex Table with Merged Cells}
\label{tab:complex}
\begin{tabular}{|l|c|c|c|}
\hline
\rowcolor{secondary!30}
\textbf{Category} & \multicolumn{3}{c|}{\textbf{Measurements}} \\
\cline{2-4}
\rowcolor{secondary!30}
& \textbf{Min} & \textbf{Max} & \textbf{Avg} \\
\hline
\multirow{2}{*}{Group A} & 10 & 50 & 30 \\
\cline{2-4}
& 15 & 55 & 35 \\
\hline
\multirow{2}{*}{Group B} & 20 & 60 & 40 \\
\cline{2-4}
& 25 & 65 & 45 \\
\hline
\end{tabular}
\end{table}

\subsection{Wide Table with Tabularx}
\begin{table}[H]
\centering
\caption{Auto-width Table Using Tabularx}
\label{tab:wide}
\begin{tabularx}{\textwidth}{|X|X|X|}
\hline
\rowcolor{accent!20}
\textbf{Long Column Header 1} & \textbf{Long Column Header 2} & \textbf{Long Column Header 3} \\
\hline
This is a long text that will automatically wrap to multiple lines within the cell & 
Another long description that demonstrates the automatic text wrapping feature & 
Final column with extensive content that spans multiple lines \\
\hline
Short text & Medium length text & More content here \\
\hline
\end{tabularx}
\end{table}

\subsection{Long Table (Multi-page)}
\begin{longtable}{lcc}
\caption{Long Table That Spans Multiple Pages} \\
\label{tab:long} \\
\toprule
\textbf{ID} & \textbf{Description} & \textbf{Value} \\
\midrule
\endfirsthead

\multicolumn{3}{c}{{\bfseries \tablename\ \thetable{} -- continued from previous page}} \\
\toprule
\textbf{ID} & \textbf{Description} & \textbf{Value} \\
\midrule
\endhead

\midrule
\multicolumn{3}{r}{{Continued on next page}} \\
\endfoot

\bottomrule
\endlastfoot

1 & Item description 1 & 100 \\
2 & Item description 2 & 200 \\
3 & Item description 3 & 300 \\
4 & Item description 4 & 400 \\
5 & Item description 5 & 500 \\
6 & Item description 6 & 600 \\
7 & Item description 7 & 700 \\
8 & Item description 8 & 800 \\
9 & Item description 9 & 900 \\
10 & Item description 10 & 1000 \\
% Add more rows as needed
\end{longtable}

\subsection{Table with Mathematical Content}
\begin{table}[H]
\centering
\caption{Table with Mathematical Expressions}
\label{tab:math}
\begin{tabular}{lcc}
\toprule
\textbf{Function} & \textbf{Derivative} & \textbf{Integral} \\
\midrule
$f(x) = x^2$ & $f'(x) = 2x$ & $\int x^2 dx = \frac{x^3}{3}$ \\
$f(x) = \sin(x)$ & $f'(x) = \cos(x)$ & $\int \sin(x) dx = -\cos(x)$ \\
$f(x) = e^x$ & $f'(x) = e^x$ & $\int e^x dx = e^x$ \\
\bottomrule
\end{tabular}
\end{table}

% ============================================================================
% SECTION 2: TCOLORBOX EXAMPLES
% ============================================================================
\section{Colored Box Templates (tcolorbox)}

\subsection{Basic Colored Box}
\begin{basicbox}[Title of the Box]
This is a basic colored box with default styling. You can put any content here,
including text, equations, lists, and even other environments.
\end{basicbox}

\subsection{Theorem-style Box}
\begin{theobox}{Pythagorean Theorem}
In a right triangle, the square of the hypotenuse equals the sum of squares 
of the other two sides:
\[ a^2 + b^2 = c^2 \]
\end{theobox}

\subsection{Example Box with Shadow}
\begin{examplebox}{Important Example}
This box has a drop shadow effect and uses a yellow color scheme. It's perfect
for highlighting important examples or notes in your document.
\begin{itemize}
    \item Feature 1: Drop shadow
    \item Feature 2: Custom colors
    \item Feature 3: Rounded corners
\end{itemize}
\end{examplebox}

\subsection{Code Listing Box}
\begin{codebox}{Python Code Example}
def fibonacci(n):
    """Generate Fibonacci sequence up to n terms"""
    if n <= 0:
        return []
    elif n == 1:
        return [0]
    
    sequence = [0, 1]
    while len(sequence) < n:
        next_val = sequence[-1] + sequence[-2]
        sequence.append(next_val)
    
    return sequence

print(fibonacci(10))
\end{codebox}

\subsection{Breakable Box for Long Content}
\begin{breakablebox}{Long Content Box}
This box can break across pages seamlessly. It's ideal for long theorems, 
proofs, or extended discussions that might span multiple pages.

\lipsum[1]

The box maintains its styling even when broken across pages, with visual 
indicators showing the continuation.

\lipsum[2]
\end{breakablebox}

\subsection{Info/Warning/Danger Boxes}
\begin{infobox}{Information}
This is an information box styled in blue. Use it for general information, 
tips, or helpful notes that readers should be aware of.
\end{infobox}

\begin{warningbox}{Warning}
This is a warning box styled in orange/yellow. Use it to alert readers about 
potential issues, caveats, or important considerations.
\end{warningbox}

\begin{dangerbox}{Danger/Critical}
This is a danger box styled in red. Use it for critical warnings, errors, or 
information that requires immediate attention.
\end{dangerbox}

% ============================================================================
% SECTION 3: FIGURES AND IMAGES
% ============================================================================
\section{Figure and Image Templates}

\subsection{Basic Figure with Caption}
\begin{figure}[H]
\centering
\begin{tikzpicture}[scale=0.8]
    \draw[fill=blue!20] (0,0) rectangle (4,3);
    \node at (2,1.5) {Placeholder Image};
    \draw[->, thick] (0,0) -- (5,0) node[right] {$x$};
    \draw[->, thick] (0,0) -- (0,4) node[above] {$y$};
\end{tikzpicture}
\caption{A simple TikZ figure with axes}
\label{fig:basic}
\end{figure}

\subsection{Subfigures}
\begin{figure}[H]
\centering
\begin{subfigure}[b]{0.45\textwidth}
\centering
\begin{tikzpicture}
    \draw[fill=red!20] (0,0) circle (1.5);
    \node at (0,0) {Figure A};
\end{tikzpicture}
\caption{First subfigure}
\label{fig:sub1}
\end{subfigure}
\hfill
\begin{subfigure}[b]{0.45\textwidth}
\centering
\begin{tikzpicture}
    \draw[fill=green!20] (0,0) rectangle (3,2);
    \node at (1.5,1) {Figure B};
\end{tikzpicture}
\caption{Second subfigure}
\label{fig:sub2}
\end{subfigure}
\caption{Multiple subfigures in one figure environment}
\label{fig:combined}
\end{figure}

\subsection{Wrapped Figure}
\begin{wrapfigure}{r}{0.5\textwidth}
\centering
\begin{tikzpicture}
    \draw[fill=orange!20] (0,0) -- (2,0) -- (2,2) -- (0,2) -- cycle;
    \node at (1,1) {Wrapped};
\end{tikzpicture}
\caption{A wrapped figure}
\label{fig:wrapped}
\end{wrapfigure}

\lipsum[3]

This text wraps around the figure on the right side. Wrapped figures are useful
when you want to integrate images more naturally into the text flow.

% ============================================================================
% SECTION 4: WORKFLOW DIAGRAMS
% ============================================================================
\section{Workflow and Process Diagrams}

\subsection{Simple Flowchart}
\begin{figure}[H]
\centering
\begin{tikzpicture}[node distance=2cm, auto]
    \node [startstop] (start) {Start};
    \node [process, below of=start] (process1) {Process 1};
    \node [decision, below of=process1] (decision) {Decision?};
    \node [process, below of=decision, yshift=-1cm] (process2) {Process 2};
    \node [startstop, below of=process2] (end) {End};
    \node [process, right of=decision, xshift=3cm] (process3) {Process 3};
    
    \draw [arrow] (start) -- (process1);
    \draw [arrow] (process1) -- (decision);
    \draw [arrow] (decision) -- node[right] {Yes} (process2);
    \draw [arrow] (decision) -- node[above] {No} (process3);
    \draw [arrow] (process3) |- (process1);
    \draw [arrow] (process2) -- (end);
\end{tikzpicture}
\caption{Simple workflow flowchart}
\label{fig:flowchart}
\end{figure}

\subsection{Horizontal Process Flow}
\begin{figure}[H]
\centering
\begin{tikzpicture}[
    node distance=2.5cm,
    block/.style={
        rectangle,
        draw,
        fill=blue!10,
        text width=5em,
        text centered,
        rounded corners,
        minimum height=3em
    },
    arrow/.style={
        thick,
        ->,
        >=stealth
    }
]
    \node [block] (init) {Initialize};
    \node [block, right of=init] (identify) {Identify};
    \node [block, right of=identify] (evaluate) {Evaluate};
    \node [block, right of=evaluate] (update) {Update};
    \node [block, right of=update] (final) {Finalize};
    
    \draw [arrow] (init) -- (identify);
    \draw [arrow] (identify) -- (evaluate);
    \draw [arrow] (evaluate) -- (update);
    \draw [arrow] (update) -- (final);
\end{tikzpicture}
\caption{Horizontal process flow diagram}
\label{fig:horizontal}
\end{figure}

\subsection{Circular Workflow}
\begin{figure}[H]
\centering
\begin{tikzpicture}
    \def\radius{2.5cm}
    \def\n{6}
    
    % Draw nodes in a circle
    \foreach \i in {1,...,\n} {
        \node[
            block,
            minimum size=2cm
        ] (node\i) at ({360/\n * (\i - 1)}:\radius) {Step \i};
    }
    
    % Draw arrows between nodes
    \foreach \i [evaluate=\i as \next using {int(mod(\i, \n) + 1)}] in {1,...,\n} {
        \draw[arrow] (node\i) -- (node\next);
    }
\end{tikzpicture}
\caption{Circular workflow diagram}
\label{fig:circular}
\end{figure}

\subsection{Swimlane Diagram}
\begin{figure}[H]
\centering
\begin{tikzpicture}[
    lane/.style={
        draw,
        minimum width=10cm,
        minimum height=2cm,
        anchor=north west
    },
    task/.style={
        draw,
        fill=blue!20,
        minimum width=2cm,
        minimum height=0.8cm,
        text centered
    }
]
    % Swimlanes
    \node[lane] (lane1) at (0,0) {Department A};
    \node[lane] (lane2) at (0,-2) {Department B};
    \node[lane] (lane3) at (0,-4) {Department C};
    
    % Tasks in lane 1
    \node[task] (t1) at (2,0) {Task 1};
    \node[task] (t4) at (6,0) {Task 4};
    
    % Tasks in lane 2
    \node[task] (t2) at (4,-2) {Task 2};
    \node[task] (t5) at (8,-2) {Task 5};
    
    % Tasks in lane 3
    \node[task] (t3) at (6,-4) {Task 3};
    \node[task] (t6) at (10,-4) {Task 6};
    
    % Arrows
    \draw[arrow] (t1) -- (t2);
    \draw[arrow] (t2) -- (t3);
    \draw[arrow] (t3) -- (t4);
    \draw[arrow] (t4) -- (t5);
    \draw[arrow] (t5) -- (t6);
\end{tikzpicture}
\caption{Swimlane workflow diagram}
\label{fig:swimlane}
\end{figure}

% ============================================================================
% SECTION 5: SCHEMAS AND ARCHITECTURE DIAGRAMS
% ============================================================================
\section{Schema and Architecture Diagrams}

\subsection{System Architecture}
\begin{figure}[H]
\centering
\begin{tikzpicture}[
    component/.style={
        rectangle,
        draw,
        fill=blue!10,
        minimum width=3cm,
        minimum height=1.5cm,
        text centered
    },
    database/.style={
        cylinder,
        shape aspect=0.5,
        draw,
        fill=green!10,
        minimum width=2cm,
        minimum height=1.5cm,
        text centered
    },
    user/.style={
        circle,
        draw,
        fill=orange!10,
        minimum size=1.5cm
    }
]
    % Nodes
    \node[user] (user) {User};
    \node[component, right of=user, xshift=2cm] (frontend) {Frontend};
    \node[component, right of=frontend, xshift=2cm] (backend) {Backend};
    \node[database, below of=backend, yshift=-1cm] (db) {Database};
    \node[component, above of=backend, yshift=1cm] (api) {API Gateway};
    
    % Connections
    \draw[arrow] (user) -- (frontend);
    \draw[arrow] (frontend) -- (backend);
    \draw[arrow] (backend) -- (db);
    \draw[arrow] (backend) -- (api);
    \draw[arrow] (api) -- (backend);
\end{tikzpicture}
\caption{System architecture diagram}
\label{fig:architecture}
\end{figure}

\subsection{Entity Relationship Diagram}
\begin{figure}[H]
\centering
\begin{tikzpicture}[
    entity/.style={
        rectangle,
        draw,
        minimum width=3cm,
        minimum height=1cm,
        text centered
    },
    relationship/.style={
        diamond,
        draw,
        fill=yellow!20,
        minimum width=1.5cm,
        minimum height=1.5cm,
        text centered
    }
]
    \node[entity] (student) {Student};
    \node[relationship, right of=student, xshift=3cm] (enrolls) {Enrolls};
    \node[entity, right of=enrolls, xshift=3cm] {Course};
    
    \node[entity, below of=student, yshift=-2cm] (professor) {Professor};
    \node[relationship, right of=professor, xshift=3cm] (teaches) {Teaches};
    
    \draw[arrow] (student) -- (enrolls);
    \draw[arrow] (enrolls) -- (3.5,0) node[right] {Course};
    \draw[arrow] (professor) -- (teaches);
    \draw[arrow] (teaches) -- (3.5,-2) node[right] {Course};
\end{tikzpicture}
\caption{Entity Relationship Diagram}
\label{fig:erd}
\end{figure}

\subsection{Network Topology}
\begin{figure}[H]
\centering
\begin{tikzpicture}[
    server/.style={
        rectangle,
        draw,
        fill=blue!20,
        minimum width=1.5cm,
        minimum height=2cm,
        text centered
    },
    router/.style={
        circle,
        draw,
        fill=green!20,
        minimum size=1cm
    },
    client/.style={
        circle,
        draw,
        fill=orange!20,
        minimum size=0.8cm
    }
]
    % Core router
    \node[router] (core) {Core};
    
    % Distribution routers
    \foreach \i in {1,2,3} {
        \node[router, below left of=core, xshift=-2cm*\i, yshift=-2cm] (dist\i) {D\i};
        \draw[thick] (core) -- (dist\i);
    }
    
    % Servers
    \node[server, above of=core, yshift=1.5cm] (srv1) {Server 1};
    \node[server, right of=core, xshift=2cm] (srv2) {Server 2};
    \draw[thick] (core) -- (srv1);
    \draw[thick] (core) -- (srv2);
    
    % Clients
    \foreach \i in {1,2,3} {
        \node[client, below of=dist\i, yshift=-1cm] (cli\i) {C\i};
        \draw[thick] (dist\i) -- (cli\i);
    }
\end{tikzpicture}
\caption{Network topology diagram}
\label{fig:network}
\end{figure}

% ============================================================================
% SECTION 6: ADVANCED TIKZ DIAGRAMS
% ============================================================================
\section{Advanced TikZ Diagrams}

\subsection{Mind Map}
\begin{figure}[H]
\centering
\begin{tikzpicture}[
    mindmap,
    concept color=blue,
    text=white,
    level 1 concept/.append style={
        sibling angle=120,
        font=\small
    },
    level 2 concept/.append style={
        sibling angle=60,
        font=\footnotesize
    }
]
    \node[concept] {Main Topic}
        child {
            node[concept] {Subtopic 1}
            child { node[concept] {Detail 1.1} }
            child { node[concept] {Detail 1.2} }
        }
        child {
            node[concept] {Subtopic 2}
            child { node[concept] {Detail 2.1} }
            child { node[concept] {Detail 2.2} }
        }
        child {
            node[concept] {Subtopic 3}
            child { node[concept] {Detail 3.1} }
            child { node[concept] {Detail 3.2} }
        };
\end{tikzpicture}
\caption{Mind map diagram}
\label{fig:mindmap}
\end{figure}

\subsection{Timeline}
\begin{figure}[H]
\centering
\begin{tikzpicture}[
    timeline/.style={
        line width=2pt,
        blue
    },
    event/.style={
        circle,
        fill=blue,
        inner sep=3pt
    },
    label/.style={
        text width=2cm,
        align=center,
        font=\small
    }
]
    % Timeline
    \draw[timeline] (0,0) -- (10,0);
    
    % Events
    \foreach \x/\year/\desc in {
        0/2020/Start,
        2/2021/Phase 1,
        5/2022/Phase 2,
        8/2023/Phase 3,
        10/2024/Complete
    } {
        \node[event] at (\x,0) {};
        \node[label, below] at (\x,-0.5) {\year};
        \node[label, above] at (\x,0.5) {\desc};
    }
\end{tikzpicture}
\caption{Project timeline}
\label{fig:timeline}
\end{figure}

\subsection{Gantt Chart}
\begin{figure}[H]
\centering
\begin{ganttchart}[
    hgrid,
    vgrid,
    x unit=1cm,
    y unit chart=0.7cm,
    time slot format=isodate,
    time slot unit=week,
    canvas/.style={
        draw=none
    },
    title/.style={
        fill=blue!20,
        font=\bfseries
    },
    bar/.style={
        fill=blue!50,
        rounded corners
    },
    bar incomplete/.style={
        fill=orange!50
    },
    bar milestone/.style={
        fill=red,
        shape=diamond
    }
]{2024-01-01}{2024-12-31}
    \gantttitle{2024}{12} \\
    \gantttitlelist{1,...,12}{1} \\
    \ganttgroup{Planning}{2024-01-01}{2024-02-28} \\
    \ganttbar{Research}{2024-01-01}{2024-01-31} \\
    \ganttbar{Design}{2024-02-01}{2024-02-28} \\
    \ganttgroup{Development}{2024-03-01}{2024-08-31} \\
    \ganttbar{Implementation}{2024-03-01}{2024-06-30} \\
    \ganttbar{Testing}{2024-07-01}{2024-08-31} \\
    \ganttgroup{Deployment}{2024-09-01}{2024-12-31} \\
    \ganttbar{Launch}{2024-09-01}{2024-09-30} \\
    \ganttmilestone{Review}{2024-10-15} \\
    \ganttbar{Optimization}{2024-11-01}{2024-12-31}
\end{ganttchart}
\caption{Project Gantt chart}
\label{fig:gantt}
\end{figure}

\subsection{Petri Net}
\begin{figure}[H]
\centering
\begin{tikzpicture}[
    place/.style={
        circle,
        draw,
        minimum size=1cm,
        fill=blue!10
    },
    transition/.style={
        rectangle,
        draw,
        minimum width=0.5cm,
        minimum height=1cm,
        fill=green!10
    },
    token/.style={
        circle,
        fill=black,
        minimum size=0.3cm
    }
]
    % Places
    \node[place] (p1) at (0,0) {P1};
    \node[place] (p2) at (4,0) {P2};
    \node[place] (p3) at (0,-3) {P3};
    \node[place] (p4) at (4,-3) {P4};
    
    % Transitions
    \node[transition] (t1) at (2,0) {T1};
    \node[transition] (t2) at (2,-3) {T2};
    
    % Arcs
    \draw[arrow] (p1) -- (t1);
    \draw[arrow] (t1) -- (p2);
    \draw[arrow] (p2) |- (t2);
    \draw[arrow] (t2) -- (p3);
    \draw[arrow] (p3) -| (t1);
    \draw[arrow] (t2) -- (p4);
    
    % Tokens
    \node[token] at (p1.center) {};
    \node[token] at (p1.center) [xshift=0.2cm] {};
\end{tikzpicture}
\caption{Petri net diagram}
\label{fig:petri}
\end{figure}

% ============================================================================
% SECTION 7: MATHEMATICAL DIAGRAMS
% ============================================================================
\section{Mathematical Diagrams}

\subsection{Commutative Diagram}
\begin{figure}[H]
\centering
\begin{tikzcd}[
    row sep=2cm,
    column sep=2cm,
    arrows={-latex}
]
    A \arrow[r, "f"] \arrow[d, "g"'] & B \arrow[d, "h"] \\
    C \arrow[r, "k"'] & D
\end{tikzcd}
\caption{Commutative diagram}
\label{fig:commutative}
\end{figure}

\subsection{Hasse Diagram}
\begin{figure}[H]
\centering
\begin{tikzpicture}[
    node distance=1.5cm,
    every node/.style={
        circle,
        draw,
        fill=white,
        minimum size=0.8cm
    }
]
    \node (top) {1};
    \node[below left of=top] (a) {a};
    \node[below right of=top] (b) {b};
    \node[below of=a, xshift=1.5cm] (c) {c};
    \node[below of=c] (bottom) {0};
    
    \draw (bottom) -- (a);
    \draw (bottom) -- (b);
    \draw (bottom) -- (c);
    \draw (a) -- (top);
    \draw (b) -- (top);
    \draw (c) -- (top);
\end{tikzpicture}
\caption{Hasse diagram}
\label{fig:hasse}
\end{figure}

% ============================================================================
% SECTION 8: CHEMICAL AND MOLECULAR DIAGRAMS
% ============================================================================
\section{Chemical Structures}

\subsection{Simple Molecule}
\begin{figure}[H]
\centering
\begin{tikzpicture}
    \node[draw, regular polygon, regular polygon sides=6, minimum size=3cm] (hex) {};
    \foreach \i in {1,...,6} {
        \fill (hex.corner \i) circle (2pt);
    }
    \node at (0,0) {Benzene};
\end{tikzpicture}
\caption{Hexagonal molecular structure}
\label{fig:molecule}
\end{figure}

% ============================================================================
% SECTION 9: ALGORITHMS AND PSEUDOCODE
% ============================================================================
\section{Algorithm Templates}

\subsection{Algorithm using algorithm2e style}
\begin{codebox}{Binary Search Algorithm}
function binary_search(arr, target):
    low = 0
    high = length(arr) - 1
    
    while low <= high:
        mid = (low + high) // 2
        
        if arr[mid] == target:
            return mid
        elif arr[mid] < target:
            low = mid + 1
        else:
            high = mid - 1
    
    return -1
\end{codebox}

% ============================================================================
% SECTION 10: ORGANIZATIONAL CHARTS
% ============================================================================
\section{Organizational Charts}

\subsection{Tree Structure Org Chart}
\begin{figure}[H]
\centering
\begin{tikzpicture}[
    level distance=1.5cm,
    level 1/.style={sibling distance=4cm},
    level 2/.style={sibling distance=2cm},
    every node/.style={
        rectangle,
        draw,
        fill=blue!10,
        minimum width=2cm,
        minimum height=0.8cm,
        text centered
    }
]
    \node {CEO}
        child {
            node {CTO}
            child { node {Dev Team} }
            child { node {QA Team} }
        }
        child {
            node {CFO}
            child { node {Accounting} }
            child { node {Finance} }
        }
        child {
            node {COO}
            child { node {Operations} }
            child { node {HR} }
        };
\end{tikzpicture}
\caption{Organizational chart}
\label{fig:orgchart}
\end{figure}

% ============================================================================
% SECTION 11: STATE MACHINES
% ============================================================================
\section{State Machines and Automata}

\subsection{Finite State Machine}
\begin{figure}[H]
\centering
\begin{tikzpicture}[
    state/.style={
        circle,
        draw,
        minimum size=1.2cm,
        fill=blue!10
    },
    initial/.style={
        initial,
        initial text=
    }
]
    \node[state, initial] (q0) {q$_0$};
    \node[state, right of=q0] (q1) {q$_1$};
    \node[state, right of=q1, accepting] (q2) {q$_2$};
    
    \path[->]
        (q0) edge[bend left] node[above] {a} (q1)
        (q1) edge[bend left] node[below] {b} (q0)
        (q1) edge[bend left] node[above] {a} (q2)
        (q2) edge[loop above] node {a,b} (q2);
\end{tikzpicture}
\caption{Finite state machine}
\label{fig:fsm}
\end{figure}

% ============================================================================
% SECTION 12: COMPARISON TABLES
% ============================================================================
\section{Comparison Tables}

\subsection{Feature Comparison Matrix}
\begin{table}[H]
\centering
\caption{Feature Comparison Matrix}
\label{tab:comparison}
\begin{tabular}{lcccc}
\toprule
\rowcolor{primary!20}
\textbf{Feature} & \textbf{Option A} & \textbf{Option B} & \textbf{Option C} & \textbf{Option D} \\
\midrule
Performance & \cellcolor{green!20}High & \cellcolor{yellow!20}Medium & \cellcolor{red!20}Low & \cellcolor{green!20}High \\
Cost & \cellcolor{red!20}\$ \$ \$ & \cellcolor{yellow!20}\$ \$ & \cellcolor{green!20}\$ & \cellcolor{yellow!20}\$ \$ \\
Scalability & \cellcolor{green!20}Excellent & \cellcolor{green!20}Good & \cellcolor{yellow!20}Fair & \cellcolor{green!20}Excellent \\
Support & \cellcolor{green!20}24/7 & \cellcolor{yellow!20}Business hrs & \cellcolor{red!20}Limited & \cellcolor{green!20}24/7 \\
Ease of Use & \cellcolor{yellow!20}Medium & \cellcolor{green!20}Easy & \cellcolor{green!20}Very Easy & \cellcolor{yellow!20}Medium \\
\bottomrule
\end{tabular}
\end{table}

% ============================================================================
% SECTION 13: CUSTOM ENVIRONMENTS
% ============================================================================
\section{Custom Environment Examples}

\subsection{Proof Environment}
\begin{breakablebox}{Proof}
\textbf{Theorem:} The sum of two even numbers is even.

\textbf{Proof:} Let $a$ and $b$ be even numbers. Then $a = 2m$ and $b = 2n$ 
for some integers $m$ and $n$. Therefore:
\[ a + b = 2m + 2n = 2(m + n) \]
Since $m + n$ is an integer, $a + b$ is divisible by 2, hence even. \qed
\end{breakablebox}

% ============================================================================
% SECTION 14: REFERENCES AND CITATIONS
% ============================================================================
\section{Cross-referencing Examples}

As shown in \cref{tab:simple}, tables can be easily referenced. 
The flowchart in \cref{fig:flowchart} demonstrates process flows.
Equations like $E = mc^2$ can be inline or displayed:
\begin{equation}
    \int_{-\infty}^{\infty} e^{-x^2} dx = \sqrt{\pi}
    \label{eq:gaussian}
\end{equation}
See \cref{eq:gaussian} for the Gaussian integral.

% ============================================================================
% CONCLUSION
% ============================================================================
\section{Conclusion}

This template collection provides a comprehensive set of LaTeX examples for:
\begin{itemize}
    \item Professional tables with various styling options
    \item Colored boxes using tcolorbox for highlights and annotations
    \item Figures and images with proper captions and subfigures
    \item Workflow diagrams and process flows
    \item System architecture and schema diagrams
    \item Advanced TikZ graphics including mind maps and timelines
    \item Mathematical diagrams and commutative diagrams
    \item Organizational charts and tree structures
    \item State machines and automata
    \item Code listings with syntax highlighting
\end{itemize}

All templates are customizable and can be adapted to your specific document 
requirements. Mix and match these examples to create professional-quality 
documents.

% ============================================================================
% APPENDIX
% ============================================================================
\appendix
\section{Package Reference}

\begin{table}[H]
\centering
\caption{Essential LaTeX Packages Used}
\label{tab:packages}
\begin{tabular}{ll}
\toprule
\textbf{Package} & \textbf{Purpose} \\
\midrule
\texttt{tikz} & Creating vector graphics \\
\texttt{tcolorbox} & Colored text boxes \\
\texttt{booktabs} & Professional tables \\
\texttt{graphicx} & Including images \\
\texttt{listings}/\texttt{minted} & Code listings \\
\texttt{caption}/\texttt{subcaption} & Figure captions \\
\texttt{pgfgantt} & Gantt charts \\
\texttt{tikz-cd} & Commutative diagrams \\
\texttt{smartdiagram} & Quick diagrams \\
\texttt{hyperref} & Hyperlinks \\
\texttt{cleveref} & Smart cross-references \\
\bottomrule
\end{tabular}
\end{table}

% ============================================================================
% BIBLIOGRAPHY (Example)
% ============================================================================
% \bibliographystyle{plain}
% \bibliography{references}

\end{document}
